\documentclass[11pt]{article}
\usepackage{amsmath,amssymb,amsthm}
\usepackage{booktabs}
\usepackage[margin=1in]{geometry}
\usepackage{enumitem}
\usepackage{hyperref}
\hypersetup{colorlinks=true,linkcolor=blue,citecolor=blue,urlcolor=blue}

\newcommand{\DI}{\Delta I}
\newcommand{\Tmin}{T_{\min}}
\newcommand{\tier}[1]{\textsf{#1}}

\title{\textbf{The Elimination Ledger}\\[2pt]
\large A Catalogue of Falsifications, Refractions, and Survivals\\
in the ACS / AISO Program}
\author{Bradley Wallace \\ \small TensorRent}
\date{June 2026 \quad\textbar\quad SIP License v1.1}

\begin{document}
\maketitle

\begin{abstract}
This paper reports negative results as first-class outputs. Across a single
elimination campaign we subjected the load-bearing claims of the ACS (Asymmetric
Codependent Systems) framework and the AISO/TENT computational stack to
pre-registered kill tests, each carrying an explicit falsification criterion fixed
\emph{before} the test was run. We adopt the discipline of \emph{adversarial
compression}: a claim is admitted as a theorem only after it survives an attempt to
destroy it, and a claim that fails is recorded --- with its killing computation --- as
a result of equal standing. The campaign retired the program's two highest--collapse
claims, both of which were the program's own: the information--asymmetry / central--charge
identity $\DI \equiv c$, and the conditional--uniformity height floor
$\Tmin=(2\pi e)^d/q$. Several numerical constants long carried as ``forced'' were shown
to be \emph{refractions} --- artifacts of a normalization or counting convention --- while
the \emph{relations} among them survived. We also gather the program's accumulated dead
ends from prior work into a single catalogue. The recurring lesson is that the
quantities which survive a change of instrument are relations, ratios, and equalities;
bare convention--laden numbers are not. Every verdict is labelled by a four--tier
verification hierarchy, and every kill test states what it does and does not claim.
\end{abstract}

\tableofcontents

\section{Introduction: overclaiming as the primary failure mode}

In a long-running theoretical program the dominant risk is not error but
\emph{overclaiming}: dressing a conjecture in the language of a theorem, or carrying a
convention-dependent number as though it were a forced prediction. The remedy adopted
here is to make destruction cheap and routine. Rather than accumulate positive support
for favored claims, we attempt to kill them, and we treat a clean kill as a success.
This is the strip-mining stance: one does not prove the gold is present; one removes
everything that is not gold, and what remains earns credibility by the elimination of
alternatives.

The campaign reported here ran a queue of targets, each a load-bearing claim of the ACS
framework or the AISO/TENT stack, and each equipped with a pre-registered kill
condition. The targets were ordered by expected collapse of candidate space per unit
cost, with deliberate weight given to claims that were \emph{our own}. A kill aimed at
one's own load-bearing claim is the highest-integrity move available, and --- as the
results below show --- the most productive: the two largest results of the campaign are
both falsifications of claims the program had previously advanced.

\section{Methodology}

\subsection{Adversarial compression}
Every claim passes through the cycle
\[
\text{conjecture} \;\longrightarrow\; \text{pre-registered test (with kill condition)}
\;\longrightarrow\;
\begin{cases}
\text{survives} \Rightarrow \text{compress to theorem (with proof)}\\
\text{fails} \Rightarrow \text{record as negative result (with the killing computation)}\\
\text{partial} \Rightarrow \text{record as observation (with explicit scope boundary)}.
\end{cases}
\]
The rules are non-negotiable: a conjecture is never described in theorem language until
its proof is complete; a negative result is documented with its conjecture, its killing
computation, the mechanism of failure, and the boundary it establishes; every scope
boundary states what the result does and does not claim.

\subsection{The four-tier verification hierarchy}
Each verdict is labelled by the strongest tier it has earned, and a lower tier is never
allowed to pass as a higher one.
\begin{center}
\begin{tabular}{cll}
\toprule
Tier & Standard & Example in this paper\\
\midrule
\tier{T1} & machine-measured, reproducible & $g_4$ moves $4/3\!\to\!2/3$ under rescaling\\
\tier{T2} & proven, or true by inspection & $\Lambda$ supported on prime powers\\
\tier{T3} & numerical, not yet theorem-level & GUE level-repulsion $\beta\approx 2.12$\\
\tier{T4} & explicitly falsified by computation & $\Tmin=(2\pi e)^d/q$ at $d=2$\\
\bottomrule
\end{tabular}
\end{center}

\subsection{The tomographic invariance engine and its kill-criterion}
A central instrument is a \emph{tomographic invariance engine}: a battery that varies
the measuring instrument (representation, normalization, convention, scale, analysis
window, taper) and asks which quantities move. The operative definition:
\begin{quote}
A quantity is a \textbf{refraction} if its value moves under a legitimate change of
instrument. It is an \textbf{invariant} (candidate) if it does not.
\end{quote}
The criterion is applied either numerically --- by recomputing under an instrument swap,
yielding a \tier{T1} verdict --- or structurally, by argument from a quantity's
construction when no recomputable pipeline is reachable, yielding a \tier{T2} verdict.
Before trusting the engine on genuinely uncertain calls, it was calibrated against known
invariants and known refractions, with negative controls and decoy targets included so
that the engine could not become its own refraction.

\subsection{Pre-registration, controls, decoys}
Each kill test fixes its falsification criterion before execution. Negative controls
(quantities that should show nothing) and decoys (near-miss targets that should not
register) accompany every numerical test. Two further disciplines were enforced
throughout: \emph{measured, never projected} --- a complexity argument or extrapolation
is not a measurement and is never reported as one --- and \emph{validate the instrument
before trusting it}, e.g.\ a new numerical method must reproduce an independent method on
a shared case before it is used at scale.

\section{Results I: the invariant/refraction classification of framework constants}
\label{sec:constants}

The ACS parameter ledger carried four quantities as algebraically ``locked.'' We turned
the engine inward and recomputed each under a legitimate change of representation or
counting convention. The verdicts (all \tier{T1}, machine-confirmed by re-running the
derivations under instrument swaps) are summarized in Table~\ref{tab:constants}.

\begin{table}[h]
\centering
\begin{tabular}{lllc}
\toprule
Quantity & As carried & Under instrument swap & Verdict\\
\midrule
$g_4=g_L=g_R$ & $4/3$ & $\to 2/3$ under $\mathrm{Tr}(T^aT^b)=\tfrac12\delta\;$vs$\;\delta$ & refraction (value)\\
$\gamma$ (Barbero--Immirzi) & $0.274067$ & $\to 0.190206$ under SU(2)$\to$SO(3) counting & refraction\\
$\lambda_\phi$ (Higgs quartic) & $2\sqrt3/27\approx0.1283$ & scales with Killing-form norm.\ & refraction\\
$\tilde h/h$ (Yukawa ratio) & $2/3$ & $d/d(\ln\mu)\approx 0$, normalization cancels & invariant\\
\bottomrule
\end{tabular}
\caption{Framework constants under the tomographic kill-criterion. Three bare values are
refractions; one dimensionless ratio is an invariant.}
\label{tab:constants}
\end{table}

\paragraph{$g_4=4/3$.} A bare gauge-coupling \emph{value} is generator-normalization
dependent; rescaling the trace form $\mathrm{Tr}(T^aT^b)=k\,\delta^{ab}$ sends $4/3\to
2/3$. The \emph{equality} $g_4=g_L=g_R$, however, is an equality of three couplings in
one shared normalization and is preserved under any consistent rescaling. The value is a
refraction; the unification relation is the invariant.

\paragraph{$\gamma=0.274$.} The value is counting-prescription dependent. Under the
SU(2) (half-integer) counting it is $0.274067$; under SO(3) (integer) counting it is
$0.190206$, a shift of $|\Delta\gamma|=0.0839$. The value moves under a change of the
state-counting prescription; only the existence of a fixed $\gamma$ \emph{given} a
prescription is structural.

\paragraph{$\lambda_\phi=2\sqrt3/27$.} The numerical value scales with the
generator/Killing-form normalization. The Koide-geometric \emph{origin} of the quartic
is the candidate invariant; the specific number, and its identification with the
physical running quartic (which crosses experiment only near $\mu\sim132$ GeV), are
scale- and convention-relative.

\paragraph{$\tilde h/h=2/3$.} A dimensionless ratio in which multiplicative normalization
cancels, additionally protected as a renormalization-group invariant
($d/d(\ln\mu)(\tilde h/h)\approx0$ on integration of the Yukawa system). This is the
strongest invariant candidate of the four and the only one that survives.

\paragraph{The pattern.} Across all four, the relations survive and the bare numbers do
not: the equality $g_4=g_L=g_R$, the ratio $\tilde h/h$, and the geometric origin of
$\lambda_\phi$ are invariant or candidate-invariant; the convention- and scale-laden
numbers $4/3$, $0.274$, $0.1283$ are refractions. This is the expected signature of
physics, in which observables are relations.

\section{Results II: falsifications}

\subsection{The height-floor scaling law $\Tmin=(2\pi e)^d/q$ (T4)}
\label{sec:q3}

The program's conditional-uniformity height floor was stated as
$\Tmin(d,q)=(2\pi e)^d/q$, with the mechanism that $\Tmin$ is the height at which the
Riemann--von Mangoldt count
\begin{equation}
N(T)\sim \frac{dT}{2\pi}\,\log\!\frac{qT}{(2\pi e)^d}
\label{eq:fwcount}
\end{equation}
becomes asymptotically valid, i.e.\ $qT>(2\pi e)^d$, below which ``too few zeros exist
to resolve the $d$-fold denser prime spectrum.''

For the Riemann zeta function ($d=1,q=1$) the floor is $2\pi e\approx 17.08$, exactly
the zero of the leading count term $(T/2\pi)(\log(T/2\pi)-1)$ --- a window-free analytic
identity, with decoy heights ($2\pi e\cdot\varphi$, $(2\pi e)^2/10$) not privileged. At
$d=1$ this floor is a genuine invariant. But $d=1$ is the degenerate case: the standard
analytic-conductor zero count is
\begin{equation}
N(T)\sim \frac{dT}{2\pi}\,\log\!\frac{q^{1/d}T}{2\pi e},
\qquad
\Tmin^{\text{std}} = 2\pi e\, q^{-1/d},
\label{eq:stdcount}
\end{equation}
and Eqs.~\eqref{eq:fwcount} and \eqref{eq:stdcount} coincide \emph{only} at $d=1$, where
both give $2\pi e/q$.

\paragraph{The discriminating test.} The $d>1$ case decides between them. We computed the
zeros of two Dirichlet $L$-functions (via Hurwitz-zeta construction and Hardy-$Z$ sign
changes) and of a genuine degree-two $L$-function, the Dedekind zeta of $\mathbb{Q}(i)$,
$\zeta_K=\zeta\cdot L(\chi_{-4})$, conductor $q=4$. The zeros of $\zeta_K$ are the union
of the $\zeta$ zeros and the $L(\chi_{-4})$ zeros, so its counting function is exact.

\begin{table}[h]
\centering
\begin{tabular}{lccccc}
\toprule
$L$-function & $d$ & $q$ & framework $(2\pi e)^d/q$ & standard $2\pi e\,q^{-1/d}$ & $N_{\text{actual}}$ at framework floor\\
\midrule
$L(\chi_{-4})$ & 1 & 4 & $4.27$ & $4.27$ & $0$ (first zero $6.02$)\\
$L(\chi_{3})$ & 1 & 3 & $5.69$ & $5.69$ & $0$ (first zero $8.04$)\\
$\zeta_{\mathbb{Q}(i)}$ & \textbf{2} & \textbf{4} & $\mathbf{72.93}$ & $\mathbf{8.54}$ & $\mathbf{51}$\\
\bottomrule
\end{tabular}
\caption{The degree-two case falsifies the scaling law. The framework floor sits above
51 actual zeros; the standard floor sits just after the first zero.}
\label{tab:q3}
\end{table}

At $d=2$ the framework floor $(2\pi e)^2/4=72.93$ sits \emph{above $51$ actual zeros} ---
flatly contradicting the framework's own definition of $\Tmin$ as the height below which
``too few zeros exist.'' Equation~\eqref{eq:fwcount} predicts $\approx 0$ zeros up to
$T=72.93$; the truth is $51$, and Eq.~\eqref{eq:stdcount} correctly predicts $\approx
50$. The standard floor $8.54$ sits just after the first zero ($N=1$), as a floor should.

\paragraph{Verdict (T4).} The floor \emph{value} at $d=1,q=1$ remains an analytic
invariant. The \emph{scaling law} $\Tmin=(2\pi e)^d/q$ is falsified in its degree
dependence: the floor scales as $2\pi e\,q^{-1/d}$, overshooting by a factor $\sim 8.5$
at $d=2$. The $\zeta$ ``survival'' in earlier scoped reporting was precisely the
degenerate case in which the wrong degree dependence coincides with the right one ---
exactly where the strip-mine anticipated the discriminator would lie.

\subsection{The information-asymmetry / central-charge identity $\DI\equiv c$ (T4 as stated)}
\label{sec:q4}

The program's stated load-bearing open conjecture identifies the information asymmetry
$\DI$ with the renormalization-group $c$-function. The reachable definitions are: a
routing scalar $\delta_i=0.4\,(\text{consistency})+0.3\,(\text{divergence})+
0.2\,(\text{late\_error})+0.1\,(\text{recomposition})\in[0,1]$, and a physical
transfer-entropy asymmetry
\begin{equation}
\DI = \mathrm{TE}(F\!\to\!G)-\mathrm{TE}(G\!\to\!F).
\end{equation}
We built an independent reference $c$-function (the Casini--Huerta entropic
$c$-function on a free-Majorana transverse-field Ising chain, computed by a
correlation-matrix method gate-validated against exact diagonalization to $10^{-13}$).
It recovers the Ising central charge $c=1/2$ at criticality (fitted $0.488$ at scale)
and flows monotonically to $0$ in the massive phase.

\paragraph{The identity fails (T2 $\to$ T4 as stated).} The central charge is unbounded
above ($c=1$ for a free boson, $c=N$ for $N$ free bosons, $c=26$ for the bosonic
string). Hence:
\begin{itemize}[nosep]
\item If $\DI$ is the bounded routing scalar ($\in[0,1]$), the identity is falsified for
every theory with $c>1$.
\item If $\DI$ is the transfer-entropy asymmetry, three structural mismatches block the
identity: \textbf{(i) sign} --- $c\ge0$, but a difference of transfer entropies is signed;
\textbf{(ii) fixed-point value} --- on a symmetric fixed point $\mathrm{TE}(F\!\to\!G)=
\mathrm{TE}(G\!\to\!F)$ so $\DI=0$, while $c$ equals the (nonzero) central charge;
\textbf{(iii) category} --- transfer entropy is a temporal, directed-information-flow
quantity requiring a time series, whereas $c$ is static and spatial (Zamolodchikov;
Casini--Huerta), so the two are not even defined on the same object.
\end{itemize}

\paragraph{The surviving reading is a theorem, not a conjecture.} The only reading under
which a monotone quantity ``plays the role of'' $c$ is the mutual-information / entropic
$c$-function. We confirmed numerically that this quantity is monotone non-increasing
along the flow for every mass tested. But under that reading $\DI=c$ is not a novel
conjecture: it is the Casini--Huerta entropic $c$-theorem, monotone by strong
subadditivity. The literal transfer-entropy asymmetry cannot furnish it: it is undefined
on a static, time-translation-invariant ground state and vanishes by symmetry on a
symmetric bipartition.

\paragraph{Verdict.} Every reading is accounted for: $\DI\equiv c$ is either
false/ill-defined (transfer-entropy asymmetry, or bounded routing scalar at $c>1$) or a
restatement of an established theorem (mutual-information reading). No reading is both
novel and true. The highest-collapse target is retired.

\subsection{The BRA speed-superiority claim (T4 for speed; determinism survives at T1)}
\label{sec:q2}

The BRA kernel was advanced as a faster deterministic replacement for a TensorFlow
operation, with a withdrawn projection of a $489\times$ speedup. We identified the
matched operation (a Gabor wave-packet render/energy, replacing
$\exp(-\delta t^2/2w^2)\exp(2\pi i f\,\delta t)$ in f32) and benchmarked the release-mode
kernel against it.

\begin{table}[h]
\centering
\begin{tabular}{cccc}
\toprule
$N$ & BRA f64 (ns/op) & TF f32 (ns/op) & ratio\\
\midrule
16 & 225.6 & 124.0 & $1.82\times$ slower\\
64 & 700.9 & 469.8 & $1.49\times$ slower\\
256 & 2505.5 & 1638.0 & $1.53\times$ slower\\
1024 & 8326.3 & 6074.5 & $1.37\times$ slower\\
\bottomrule
\end{tabular}
\caption{Measured BRA vs.\ TF throughput. The speed-superiority claim is killed.}
\label{tab:q2}
\end{table}

\paragraph{Verdict.} BRA is $1.37$--$1.82\times$ \emph{slower}; the $489\times$ claim is
obliterated. One confound is recorded honestly: the comparison is BRA-f64 against
TF-f32, and f64 carries $\sim1.5$--$2\times$ the work, so ``competitive at matched f64''
is untested (plausibly near parity). The genuine surviving value is \emph{determinism}:
the integer and f64 paths are bit-identical across runs (\tier{T1}). The speed claim is
dead; the determinism claim --- always the real one --- stands.

\subsection{Residual prime-orbit off-diagonal mechanisms (T2 + T3)}
\label{sec:q6}

We tested whether any \emph{peaked} off-diagonal mechanism contributes to the
zeta spectrum beyond the diagonal prime-power terms. Using the explicit-formula dual
periodogram on the zeros, the power ratios to local baseline were: fundamentals
$\log 2,3,5,7$ at $\sim 3$--$4\times10^8$; prime-power harmonics $2\log2,3\log2,2\log3,
2\log5$ at $\sim 5\times10^7$--$4\times10^8$ (real peaks, suppressed by the $p^{k/2}$
weight); composites $\log6,10,12,15$ at $0.3$--$5.4$ (baseline noise); decoys $\sim1$.

\paragraph{Verdict.} The peaked spectrum is exactly the diagonal prime-power set
$\{k\log p\}$ at $10^7$--$10^8$ contrast. There is a structural reason, not merely a
null: the explicit-formula dual carries weight only on the von Mangoldt function
$\Lambda(n)$, which is supported on prime powers, so a sum-frequency peak at
$\log(pp')$ is \emph{forbidden}. Difference-tone off-diagonal peaks were falsified in
prior work. The remaining off-diagonal content is smooth (the Montgomery pair-correlation
term, separately confirmed). No residual peaked off-diagonal mechanism exists; the tunnel
is mapped empty.

\section{Results III: class-level eliminations --- the Hilbert--Pólya operator}
\label{sec:q5}

One cannot construct and diagonalize a self-adjoint operator with spectrum equal to the
zeros --- that is the open problem itself --- but the data forces discriminating
constraints, and those constraints kill candidate \emph{classes}.

\paragraph{Symmetry class (T3).} The unfolded nearest-neighbour spacings of the $\zeta$
zeros match the GUE Wigner surmise (L2 distance $7.3\times10^{-2}$) far better than GSE
($1.7\times10^{-1}$), GOE ($2.3\times10^{-1}$), or Poisson ($6.5\times10^{-1}$); the
fitted level-repulsion exponent is $\beta\approx2.12$, with variance $0.1600$. Under
random-matrix universality the operator must lie in the unitary class with broken
time-reversal symmetry. This kills any real-symmetric / time-reversal-invariant candidate
(GOE, $\beta=1$) and any symplectic candidate (GSE, $\beta=4$).

\paragraph{Berry--Keating $xp$ (T2 + T3).} The semiclassical $xp$ count reproduces the
smooth Riemann--von Mangoldt staircase to $<0.2\%$, so it passes the smooth-density
necessary condition. But $xp$ has a continuous spectrum and no discrete eigenvalues, and
the fluctuation $S(T)$ is prime-driven, which bare $xp$ does not encode. ``$xp$ alone is
the operator'' is killed as sufficient; ``$xp$ as the smooth skeleton'' survives.

\paragraph{Verdict.} The surviving operator must simultaneously be unitary class (broken
$T$), reproduce $(T/2\pi)\log(T/2\pi e)$, and encode the primes in its fluctuation
spectrum. Two classes are mapped empty; the full construction remains open.

\section{Results IV: a tension resolved by decoupling (T3 + T2)}
\label{sec:q7}

The naive Type-I see-saw active--sterile mixing $\sin^2 2\theta\approx4\times10^{-6}$ is
excluded by the X-ray bound $\sin^2 2\theta<10^{-10}$. The proposed Pati--Salam
gauge-sector suppression is $(M_W/M_{W_R})^2$. Requiring it to clear the bound gives
$M_{W_R}\gtrsim 16$ TeV (equivalently $v_R\gtrsim 49$ TeV), comfortably below the
independent lower bounds on the Pati--Salam scale from rare decays
($K_L\to\mu e$ pushes $v_R$ to $\gtrsim$ hundreds of TeV). So the mechanism clears the
bound.

\paragraph{Verdict.} The X-ray bound is not a live obstruction once $v_R$ sits at its
independently-required scale. But the resolution is by \emph{decoupling}: clearing the
bound requires high $v_R$, hence a heavy right-handed mass $M_R\sim v_R$, so the sterile
state is no longer a keV-scale dark-matter candidate. The tension is removed for the
field-theory consistency reading, and the keV-sterile-dark-matter reading does not
survive.

\section{The accumulated dead ends}
\label{sec:deadends}

For completeness we gather the program's prior falsifications --- the empty tunnels
mapped before this campaign --- with the mechanism that killed each.

\begin{table}[h]
\centering
\small
\begin{tabular}{p{5.4cm}p{8.2cm}}
\toprule
Killed claim & Killing mechanism\\
\midrule
$\mathrm{ad}^3=2\,\mathrm{ad}$ for integers & requires $\lambda=1/\sqrt2$; eigenvalue constraint fails for integers\\
Universal $2\pi$ inversion & $\mathfrak{sl}(4,\mathbb{R})$ adjoint is hyperbolic, not rotational; representation-specific\\
Wronskian bracket $=$ Poisson bracket & Leibniz rule fails by $-fgh'$; structural claim demoted to analogy only\\
IR lattice imprint in zeros & $z$-scores null against PDG values; statistical null test\\
Intrinsic chirality of the bias & bias vanishes under covariant conjugation; not invariant\\
Signature selection via Killing form (Route A) & Cartan involution $\ne$ parity involution\\
Signature selection via stability (Route C) & all signatures give equal holonomy norms\\
Coleman--Weinberg $6\to5$ parameter reduction & $\tan\beta$ is a gauge-protected flat direction (RG-invariant ratio, multiplicative two-loop QCD, $\beta$-independent thresholds)\\
$\alpha_2$ optional & forbidden by representation theory, $T^A\Phi=0$\\
$\beta_c$ explanatory & tree-level no-go from mass-matrix rank\\
Yukawa rescue of $\tan\beta$ & impossible via the identity $M_u-M_d=(h-\tilde h)(\kappa_1-\kappa_2)$\\
Action-proximity ($\tau$) off-diagonal mechanism & inconsistent with the incommensurable surrogate null\\
Difference-tone / Tartini off-diagonal mechanism & inconsistent with the surrogate null\\
``$3$-$6$-$9$'' attractor & retired as falsified (no structural support)\\
\bottomrule
\end{tabular}
\caption{Prior dead ends gathered from the program. Each is a first-class negative
result with an explicit killing mechanism.}
\label{tab:deadends}
\end{table}

\section{Discussion: recurring patterns}

\paragraph{Relations survive; bare numbers are refractions.} The single most consistent
finding (Sec.~\ref{sec:constants}) is that the quantities surviving a change of
instrument are relations, ratios, and equalities, while convention- or scale-laden
numbers are not. This is precisely what one expects of physics, whose observables are
relations; it also gives a fast triage heuristic for which claims to suspect.

\paragraph{The degenerate-case trap.} The height-floor law (Sec.~\ref{sec:q3}) survived
every test available at $d=1$ because at $d=1$ a wrong degree dependence coincides
exactly with the right one. A claim can pass indefinitely on the degenerate slice of its
parameter space; the discriminating test lives off that slice, and naming where it lives
\emph{before} testing is what turned a scoped survival into a clean falsification.

\paragraph{Category mismatch.} The $\DI\equiv c$ identity (Sec.~\ref{sec:q4}) failed most
fundamentally on a category error: a temporal, directed quantity identified with a
static, spatial one. When an identity is proposed between objects from different
categories, the burden is a bridge construction, not a numerical coincidence.

\paragraph{Necessary but insufficient.} The Berry--Keating $xp$ operator
(Sec.~\ref{sec:q5}) passes a necessary condition (the smooth count) and is nonetheless
insufficient (no discrete spectrum, no prime fluctuations). Passing a necessary test is
not survival; a candidate must clear every condition the data forces.

\paragraph{Methodological self-corrections.} The campaign caught and discarded several of
its own artifacts before they could be read as signal: a band-relative-error onset
statistic dominated by low-height discreteness; a sub-sample indexing bug in a deviation
metric; a ``broken-lens'' control that failed to break the GUE statistic for an
identifiable reason; projected speedups mistaken for measurements; and a SIGPIPE harness
artifact that reported a passing test as a failure. Each was recorded, not hidden. The
discipline that catches these is the same one that produces the kills: pre-registration,
negative controls, and the refusal to let analysis pose as measurement.

\section{What survives, and what remains open}

\paragraph{Survivors (earned by elimination).} The dimensionless ratio $\tilde h/h=2/3$
(RG-invariant); the coupling-unification equality $g_4=g_L=g_R$; the analytic floor
$2\pi e$ at $d=1,q=1$; the Koide-geometric origin of $\lambda_\phi$; the GUE statistics
and the prime-power structure $\{k\log p\}$ of the zeta spectrum (the latter forced by
$\Lambda$-support). These survive not because they were proven from first principles
here, but because the alternatives around them were eliminated.

\paragraph{Open.} The Hilbert--Pólya operator construction (constrained but not
constructed); the action principle for the grading functional; the extension of the
spectral results to general $L$-functions (now with a corrected --- standard ---
height-floor scaling); and the conceptual question of whether a static, positive,
canonically-normalized information quantity exists that could non-trivially realize a
$c$-function role.

\section{Conclusion}

The campaign retired the two highest-collapse claims of the program --- the
$\DI\equiv c$ identity and the $\Tmin=(2\pi e)^d/q$ scaling law --- and both were the
program's own. Several ``forced'' constants were reclassified as refractions, with their
underlying relations surviving. A speed claim was killed and a determinism claim
confirmed. Two operator classes were eliminated and a peaked-spectrum tunnel mapped
empty. None of these required a new positive theory; each followed from making
destruction cheap, pre-registering the kill condition, and treating a clean negative as a
result worth the same care as a proof. The survivors are credible in proportion to what
was removed around them. That is the whole of the method: do not prove the gold is there;
remove everything that is not gold, and keep an honest ledger of every tunnel.

\appendix
\section{Reproducibility}
All numerical tests are reproducible from the accompanying scripts. The degree-two
falsification (Sec.~\ref{sec:q3}) uses \texttt{q3\_scaling\_test.py} (mpmath
Dirichlet-$L$ zeros via Hurwitz-zeta construction and Hardy-$Z$ sign changes; $\zeta$
zeros from the first $10^5$ accurate to $3\times10^{-9}$; the degree-two case built as
$\zeta\cdot L(\chi_{-4})$). The reference $c$-function (Sec.~\ref{sec:q4}) uses
\texttt{q4\_cfunction\_testbed.py} (free-fermion BdG correlation-matrix method, gate-validated
against exact diagonalization). Verdicts and tiers are recorded in the append-only
elimination ledger (\textit{The Elimination Ledger}). Random seed $20260423$ where randomness appears.

\end{document}
